\section{Ethics and Privacy}
\label{sec:ethics}

\textbf{Data.} All face images come from commercially licensed before/after
datasets whose terms cover AI processing and reproduction in this publication; no
private or non-consented patient data was used at any stage, and images sent to
hosted endpoints were transmitted inline without server-side retention options
enabled. The figures reproduce faces under that license.

\textbf{No cherry-picking.} Every experiment runs its full face $\times$ model
$\times$ procedure $\times$ control matrix under a fixed seed into an immutable,
git-stamped run directory,
and the scoring gallery renders every cell. The examples in
Figs.~\ref{fig:teaser} and \ref{fig:grid} were chosen for print legibility, but
the statistics in Table~\ref{tab:main} include every scored edit, failures and
all.

\textbf{Intended use and misuse.} A preview is an expectation-setting aid, not a
promised outcome. The clinical literature has debated morphing on exactly this
point for decades~\cite{agarwal2007morph}: an over-idealized preview can create
expectations no surgeon can meet, which is why our prompts ask for realistic
tissue repositioning rather than idealization, and why the rubric reserves a
safety-of-expectation axis. The same editing capability applied to a photo of a
non-consenting person is straightforwardly harmful; the pipeline described here
does nothing to enable that beyond what the underlying public endpoints already
permit, but deployments should verify subject consent as a precondition, and model
documentation practices~\cite{mitchell2019modelcards} apply to systems built on
these APIs just as they do to the models themselves.

\textbf{Fairness.} Because the demographic makeup of the test set is unrecorded,
we make no claim that the reported behavior holds uniformly across
populations~\cite{buolamwini2018gendershades}; a stratified test set is a
precondition for any clinical deployment.
